\scanpage{77}
La relation \deleted{\unsure{avec}} entre cette description du fibré en droites projectives $X$, et la description comme $\underline{\operatorname{Bor}}(G)$ ou $\underline{\operatorname{Bor}}(\widetilde G)$, est explicitée dans la

\uline{Proposition} Les ss-groupes de Borel \inserted{$B$} de $G$ correspondent \unsure{biunivqt|biunivoquemt|biunivoquement} par image inverse à \deleted{ceux} \inserted{$B_1$} de $\mathcal L=\mathbb V(M)^\times$ ou \inserted{$\widetilde B$} de $\widetilde G$. Un ss-groupe de Borel est connu quand on connaît son algèbre de Lie $L$ (qui est un ss fibré vectoriel de corang $1$ dans $\mathfrak g$), ou celle de $L_1$ de $\mathcal L$ qui lui correspond (qui est un ss-fibré \deleted{de rang \unsure{2|3}} \inserted{vectoriel} de corang $1$) image inverse de $L_1$ dans $M=\underline{\operatorname{Lie}}(\mathcal L)$, \deleted{ou celle} ou \uuline{en car.} $\ne2$ quand on connaît celle $\widetilde L$ de $\widetilde B$ dans $\widetilde{\mathfrak g}$, qui n'est autre que $\widetilde L\cap\widetilde{\mathfrak g}$ -- (et pour cause c'est \unsure{pareil|permis}, puisqu'alors $f:\widetilde{\mathfrak g}\xrightarrow{\sim}\mathfrak g$ !). \deleted{Pour} \inserted{En termes} \deleted{\unsure{pour un ss-fibré vectoriel}} de $\widetilde L$ on aura
\[
\widetilde B=\check{\mathbb V}(\widetilde L)\cap\underbrace{\mathbb V(M)}_{\mathcal L}^{*}\qquad(B=\widetilde B/G_m)\tag*{(57)}
\]
% archive sheet number: 76
